\section{Conclusion}
\label{sec:conclusion}

Redistricting analysis currently operates in an inferential gap:
it produces descriptive percentile claims that lack the statistical
rigour required for litigation-grade evidence.
The four challenges identified in this paper---combinatorial test statistics,
simulation-dependent null distributions, multiple testing, and small-tail
ESS---are solvable with existing statistical methodology.

The three-component framework (calibrated null distribution, FDR correction,
power analysis) bridges this gap.
It is not a new statistical theory; it is an application of well-established
methods (Benjamini-Hochberg, ESS correction, permutation tests) to the
redistricting domain.

The S-track papers (S.1--S.4) implement each component with worked examples
and provide reference implementations in \bisect{}.
Together, they convert the redistricting ensemble from a descriptive tool
into a rigorous inferential instrument---one that can support expert
testimony meeting the Daubert standard and provide courts with the
``judicially manageable standard'' that \textit{Rucho} found lacking
in partisan gerrymandering claims.

The practical recommendation for redistricting practitioners:
(1)~certify ensemble convergence using G.4's diagnostics before
reporting any percentile claim; (2)~use ESS-corrected sample sizes
for confidence intervals; (3)~apply BH correction when reporting
multiple metrics; and (4)~report power as a companion to any
null-hypothesis test.
These steps convert anecdotal evidence of extremity into statistically
defensible claims.
